\documentclass{beamer}
\usepackage{amsmath,graphicx,booktabs}
\title{Perpetual American Put}
\subtitle{A verified finite-difference obstacle solver}
\author{Coursework team}
\date{Project 2 technical draft}
\graphicspath{{../figures/}{../stage2_D_runs/20261105T132543802032Z/figures/}}
\begin{document}
\frame{\titlepage}
\begin{frame}{The decision to exercise creates an obstacle}
\begin{itemize}
\item Exercise gives $(K-S)^+$ immediately, continuation may be more valuable.
\item The stationary one-dimensional problem is $\min(-\mathcal L V,V-\psi)=0$.
\item We seek both the value and the exercise threshold.
\end{itemize}
\end{frame}
\begin{frame}{Monotone differences preserve the obstacle structure}
\begin{itemize}
\item Upwind drift produces a tridiagonal M-matrix.
\item Policy iteration alternates exercise and continuation rows.
\item The stopping test measures complementarity residual.
\item The contact node uses price-gap tolerance $10^{-7}$.
\end{itemize}
\end{frame}
\begin{frame}{Exact prices converge at roughly first order}
\begin{itemize}
\item Stage 1 exact threshold is $500/7$.
\item Four meshes reduce maximum price error from 0.05847 to 0.00736.
\item Last contact node on the finest mesh is 71.4375.
\item Nodal threshold errors need not decrease monotonically.
\end{itemize}
\end{frame}
\begin{frame}{Stage 2 converges on a fixed truncated domain}
\begin{itemize}
\item Local volatility and $q=0.02$ produce a genuinely numerical case.
\item At $L=600$, finest mesh versus fine-grid difference is 0.004604.
\item A second reference changes the price comparison by $3.28\cdot10^{-4}$.
\item Fine-reference residuals reach a reported floating-point floor.
\end{itemize}
\end{frame}
\begin{frame}{The outer boundary changes the answer materially}
\begin{center}\begin{tabular}{rrr}\toprule
$L$ & numerical $V(100)$ & collocation $V(100)$\\\midrule
300 & 15.4560 & 15.4512\\
400 & 15.8310 & 15.8260\\
600 & 16.0653 & 16.0601\\\bottomrule
\end{tabular}\end{center}
Domain enlargement and mesh refinement measure different sensitivities.
\end{frame}
\begin{frame}{Independent checks support the finite-domain result}
\begin{itemize}
\item Analytic Stage 1 solution checks price and threshold.
\item Exhaustive small-grid active sets check policy iteration.
\item Shooting and adaptive collocation agree for the same domain.
\item Shape, complementarity and tolerance diagnostics are recorded.
\end{itemize}
\end{frame}
\begin{frame}{Two solvers agree at a common residual target}
\begin{itemize}
\item Policy iteration and projected SOR use the same LCP and payoff initialisation.
\item At $h=5$ and $2.5$, maximum solution differences are below $4\cdot10^{-8}$.
\item Policy systems use compiled linear algebra, while PSOR uses Python sweeps.
\item Timings include the implementation and platform effects.
\end{itemize}
\end{frame}
\begin{frame}{Reproducibility and discussion}
\begin{itemize}
\item Reproduce with the repository README and inspect saved numerical evidence.
\item Grid and tolerance studies alone cannot remove truncation bias.
\item Explain stencil signs, complementarity and threshold extraction.
\item Complete the student audit, individual contributions and presentation rehearsal.
\end{itemize}
\end{frame}
\end{document}
